\chapter{決算発表とイベントスタディ}\label{ch:event}

この研究の中心の実験（実験 2〜5・10・11）は、決算発表という\textbf{イベント}の前後の株価の動きを調べる\term{イベントスタディ}である。この章では、米国の決算発表の仕組みと、イベントスタディの考え方と道具を説明する。

\section{米国の決算発表}

\subsection{発表の流れ}

米国の上場会社は、四半期ごとに業績を公表する。流れはおおむね次のとおりである。

\begin{enumerate}
\item \textbf{決算発表（プレスリリース）}：四半期の終わりから 3〜6 週間後。売上、利益、1 株あたり利益（EPS）と、会社によっては次の四半期や通期の見通し（\term{ガイダンス}）を発表する。通信社の配信と同時に、SEC に \textbf{8-K}（臨時報告書）の \textbf{Item 2.02}（業績の公表）として提出し、発表文そのものを添付書類 \textbf{EX-99.1} として付ける。
\item \textbf{説明会（カンファレンス・コール）}：発表の直後に、経営陣がアナリストの質問に答える。
\item \textbf{四半期報告書（10-Q）・年次報告書（10-K）}：発表から 1〜3 週間後に提出する、監査（10-K）やレビューを受けた正式な財務諸表。XBRL という機械が読める形式の数値が付く。
\end{enumerate}

\subsection{発表の時刻}

発表は、市場が開く前（\term{寄り前}、東部時間 9:30 より前）か、閉じた後（\term{引け後}、16:00 以後）に行われるのが普通である。場中に出すと、株価が急に動いて混乱するからである。

\begin{defbox}[定義：反応の日（event\_date）]
株価の終値が初めて発表を織り込む取引日を day 0 とする。
\begin{itemize}
\item 寄り前・場中に受理された発表：受理日が取引日ならその日、そうでなければ次の取引日。
\item 引け後（16:00〜17:30）・17:30 以後に受理された発表：受理日より後の最初の取引日。
\end{itemize}
SEC の受理時刻（UTC）を東部時間に直し（夏時間を考慮）、この規則で決める。
\end{defbox}

8-K の受理は、プレスリリースの配信より遅れることがある（たとえば 16:05 に配信し、翌朝 8:00 に受理）。その場合、day 0 が 1 日遅れうる。このずれを吸収するため、反応は day 0 の前日の終値から day 1 の終値までの 2 取引日で測る。

\section{驚き（サプライズ）}

株価が動くのは、発表の内容が\textbf{予想と違った}ときである。予想との差を\term{驚き}（サプライズ）と呼ぶ。

\begin{longtable}{P{3.4cm}P{6.0cm}P{5.0cm}}
\toprule
\textbf{驚き} & \textbf{定義} & \textbf{性質} \\
\midrule
\endhead
\bottomrule
\endfoot
SUE（標準化した予期しない利益） & $(\text{EPS}_q - \text{EPS}_{q-4}) / $ 過去 8 四半期の季節差の標準偏差 & 予想として「前年同期の EPS」を使う（季節の酔歩）。アナリストの予想がなくても計算できる \\
売上の驚き & 売上の前年同期比 − 過去 4 四半期の前年同期比の平均 & 売上の伸びの加速・減速 \\
アナリスト予想との差 & $(\text{EPS}_{\text{実績}} - \text{EPS}_{\text{予想}}) / |\text{EPS}_{\text{予想}}|$ & 市場の期待に最も近い。データは Alpha Vantage（実験 3） \\
\end{longtable}

\begin{notebox}[驚きの先読み]
10-Q の数値は発表の 1〜3 週間後に提出されるが、その四半期の値は発表文で公表済みなので、発表の時点の驚きに使っても先読みにはならない。ただし、SUE の分母の標準偏差に使う過去の四半期は、発表より前に提出されたものだけにする。
\end{notebox}

\section{異常リターン}

発表の効果を測るには、発表がなくても起きていたはずの動き（市場全体やセクターの動き）を除く必要がある。

\begin{equation}
\text{AR}_{i,t} = R_{i,t} - \E[R_{i,t} \mid \text{発表がない場合}]
\end{equation}
を\term{異常リターン}（アブノーマル・リターン）と呼ぶ。「発表がない場合」の見積もりには、いくつかの方法がある。

\begin{description}[leftmargin=1.5em,style=nextline]
\item[市場調整] $\text{AR} = R_i - R_m$。市場（\SP{}）のリターンを引く。
\item[市場モデル] 発表の前の期間で $R_i = a + bR_m$ を推定し、$\text{AR} = R_i - (\hat a + \hat b R_m)$。
\item[セクター調整] $\text{AR} = R_i - R_{\text{セクター ETF}}$。同じ日に同じ業種の決算が集まる効果や、業種のニュースも除ける。\textbf{このプロジェクトはこれを使った}（ETF の株価がない日は \SP{} との差）。
\end{description}

期間で足したものを\term{累積異常リターン}（CAR）と呼ぶ。窓 $[\tau_1,\tau_2]$ の CAR は $\sum_{t=\tau_1}^{\tau_2}\text{AR}_t$ である。

\section{発表後のドリフト}

1968 年の Ball と Brown 以来、良い驚きの後は株価が数週間〜数か月上がり続け、悪い驚きの後は下がり続けることが知られている。これを\term{発表後のドリフト}（PEAD）と呼ぶ。Bernard と Thomas（1989）は、SUE の上位と下位の銘柄で、発表後 60 日に数 \% の差を報告した。

ドリフトは、効率的市場仮説に対する代表的な例外で、「投資家は決算の情報に反応しきれない」ことを示すと解釈されてきた。ただし近年は、アルゴリズムの取引の増加とともに、小さくなったという報告が多い。この研究の実験 2・3 では、LLM の点数で day 2〜21 のドリフトを説明できなかった（+10bp、$t$ 0.83）。

\section{文章の分析}

数値の驚きだけでなく、発表文の\textbf{文章}も株価を動かす。経営陣の語調、ガイダンスの引き上げ・引き下げ、需要の見通しの書き方は、数値に表れない情報を含む。文章を数にする方法は第 V 部（\chapref{ch:text}）で説明する。先行研究では、
\begin{itemize}
\item Loughran と McDonald（2011）：一般の感情の辞書は金融の文章では誤る（「負債」「税」は否定語ではない）。金融用の辞書を作った。
\item Tetlock ほか（2008）：ニュースの否定語の割合が、利益と株価を予測する。
\item Lopez-Lira と Tang（2023）：ChatGPT にニュースの見出しを読ませた点数が、翌日のリターンを予測する。
\end{itemize}
などがある。

\section{イベントスタディの統計}

\subsection{横断面の回帰}

各発表を 1 行として、異常リターンを点数と統制の変数で回帰する。
\begin{equation}
\text{CAR}_i = \beta\,\text{点数}_i + \gamma\T\text{統制}_i + \text{固定効果} + \varepsilon_i
\end{equation}
統制には、数値の驚き、ほかの文章の点数、発表前のリターン、規模（時価総額の対数）を入れ、業種と年・四半期の固定効果を入れる。$\beta$ が「統制を超えて点数が説明する分」である。点数は標準化しておき、$\beta$ を「1 標準偏差あたりの bp」で報告する。

\subsection{同じ日の発表の相関}

同じ日に発表した会社どうしは、同じ市場のニュースにさらされるので、誤差が相関する（Kolari と Pynnönen 2010）。普通の標準誤差は、この相関を無視して小さく出る。そこで\textbf{発表日ごとにまとめた標準誤差}（クラスター頑健な標準誤差）を使う（\chapref{ch:stat}）。

\subsection{置換の対照}

点数を、\textbf{同じ発表日の中で}会社の間でランダムに入れ替えて回帰し直すことを 1,000 回くり返し、本物の係数がその分布のどこにあるかを見る。日付の効果は保ったまま、点数と会社の対応だけを壊すので、分布の仮定に頼らない検定になる。

\subsection{外れ値}

決算の反応には、$\pm30\%$ を超えるような極端な値がある。少数の極端な値で係数が決まらないよう、目的変数と驚きを上下 1\% で\term{丸める}（ウィンザライズ：1\% 点より外の値を 1\% 点の値に置き換える）。

\begin{projbox}
\begin{itemize}
\item \textbf{データ}：SEC EDGAR から、\SP{} の構成銘柄（期間中に一度でも入った 668 社）の 8-K Item 2.02 と EX-99.1 を取得し、受理時刻から day 0 を決めた（第 VI 部）。
\item \textbf{目的変数}：day 0 の前日の終値から day 1 の終値までの、セクター ETF に対する超過リターン（上下 1\% で丸める）。
\item \textbf{統制}：SUE、売上の驚き（欠けの印つき）、辞書・FinBERT・TF-IDF の点数、$[-21,-1]$ の超過リターン、時価総額の対数、セクターと年・四半期の固定効果。実験 3 以降はアナリスト予想との差も。
\item \textbf{判定}：点数の係数の片側 $t$（発表日でまとめた標準誤差）と、同じ発表日の中の置換 1,000 回の両方で $p<0.05$。
\item \textbf{ドリフト}：day 2〜21（実験 2・3）、day 1・day 1〜5（実験 8）。
\end{itemize}
\end{projbox}
